\begin{abstract}
The task of mathematical reasoning remains particularly difficult for language models because of its inherently intricate and formal characteristics. Here, we present DeepSeekMath~7B, an extension of DeepSeek-Coder-Base-v1.5 7B, further pre-trained with 120B math-focused tokens compiled from Common Crawl, as well as related natural language and programming data. DeepSeekMath~7B attains a noteworthy 51.7\% score on the challenging, competition-grade MATH benchmark without the aid of supplementary toolkits or ensemble voting approaches, coming close to the performance of Gemini-Ultra and GPT-4. Employing self-consistency sampling with 64 generations, DeepSeekMath~7B reaches 60.9\% accuracy on MATH. The advanced mathematical reasoning abilities of DeepSeekMath~are primarily attributable to two central innovations: First, the effective exploitation of large-scale, public web datasets using a carefully designed selection pipeline. Second, the development of Group Relative Policy Optimization (GRPO)—an adaptation of Proximal Policy Optimization (PPO)—that boosts mathematical reasoning while simultaneously improving PPO’s memory efficiency.
\end{abstract}

\section{Introduction}

Large language models (LLMs) have transformed artificial intelligence methodologies for mathematical reasoning, driving considerable progress on both quantitative reasoning \citep{MATH} and geometry reasoning \citep{trinh2024solving} evaluation sets. Additionally, LLMs have become invaluable tools for human users tackling advanced mathematical questions \citep{tao}. Nevertheless, leading-edge models like GPT-4 \citep{gpt4} and Gemini-Ultra \citep{gemini} remain closed-source, and available open-source models exhibit significantly inferior performance.

This paper presents DeepSeekMath, a language model tailored to the mathematics domain that markedly surpasses existing open-source models and closely rivals GPT-4 in terms of performance on established academic evaluations. Central to this advancement is the construction of the DeepSeekMath Corpus, an extensive and high-quality pre-training dataset consisting of 120B math tokens. The dataset is sourced from Common Crawl (CC) utilizing a classifier based on fastText \citep{joulin2016fasttext}. In its initial phase, the classifier is trained with positive samples from OpenWebMath \citep{openwebmath} and a broad assortment of unrelated webpages as negative samples. Leveraging this classifier, further positive instances are extracted from the CC, and then subject to additional curation through manual annotation. The classifier is subsequently retrained on this expanded and refined data to enhance its precision. Evaluation outcomes reveal that the resulting corpus is of notable quality; DeepSeekMath-Base 7B achieves 64.2\% accuracy on GSM8K \citep{gsm8k} and 36.2\% on the advanced-level MATH dataset \citep{MATH}, surpassing the performance of Minerva 540B \citep{minerva}. Furthermore, since the DeepSeekMath Corpus encompasses multiple languages, we observe improved results on Chinese mathematics tasks \citep{wei2023cmath,agieval}. We consider the methodologies we developed for mathematical data curation as a valuable resource for ongoing and future research, and we anticipate substantial progress in this field going forward.

DeepSeekMath-Base is initialized using DeepSeek-Coder-Base-v1.5 7B \citep{deepseek-coder}, given our observation that models pre-trained on code perform better than those derived from general-purpose LLMs for downstream tasks. Additionally, we note that mathematical pre-training not only augments the model's mathematical proficiency, but also results in gains on broader benchmarks such as MMLU \citep{mmlu} and BBH \citep{bbh}, suggesting improved general reasoning abilities.

Subsequent to pre-training, DeepSeekMath-Base undergoes mathematical instruction fine-tuning, incorporating chain-of-thought \citep{cot}, program-of-thought \citep{pot,pal}, and tool-augmented reasoning datasets \citep{tora}. The final model, DeepSeekMath-Instruct 7B, outperforms all other 7B models and achieves parity with leading 70B open-source instruction-tuned systems.

We further present Group Relative Policy Optimization (GRPO), a modified reinforcement learning approach inspired by Proximal Policy Optimization (PPO) \citep{schulman2017proximal}. Unlike PPO, GRPO omits the use of a critic network and instead estimates the baseline using group-level scores, which significantly curtails training computational requirements. Applying GRPO with only a portion of English instruction-tuning data yields remarkable performance gains for DeepSeekMath-Instruct across both in-domain (GSM8K: 82.9\% $\rightarrow$ 88.2\%, MATH: 46.8\% $\rightarrow$ 51.7\%) and out-of-domain (CMATH: 84.6\% $\rightarrow$ 88.8\%) mathematical evaluation tasks during reinforcement learning. 

Moreover, we establish a unified perspective encompassing multiple approaches—such as Rejection Sampling Fine-Tuning (RFT) \citep{yuan2023scaling}, Direct Preference Optimization (DPO) \citep{dpo}, PPO, and GRPO—framing them all as either explicit or simplified variants of reinforcement learning. Our systematic analysis involves comprehensive experiments, including comparisons of online versus offline optimization, supervision over process versus outcomes, and single-turn versus iterative reinforcement learning, to probe fundamental aspects of this framework. Finally, we elucidate how our reinforcement learning paradigm enhances instruction-tuned models, and outline future research directions aimed at developing more effective RL techniques within this unified scheme.

\subsection{Contributions}
This work presents advances in large-scale mathematical pre-training as well as an in-depth examination of reinforcement learning strategies.

\textbf{Math Pre-Training at Scale}

\begin{itemize}[topsep=0pt]
    \item We demonstrate that the widely available Common Crawl dataset encompasses substantial mathematical content. Through a rigorously developed data curation pipeline, we assemble the DeepSeekMath Corpus—a robust and high-fidelity dataset containing 120B tokens specifically filtered from mathematical web pages. This dataset substantially surpasses previous resources, being nearly sevenfold larger than the mathematical web subset employed by Minerva \citep{minerva}, and approximately nine times greater than the recently introduced OpenWebMath \citep{openwebmath}.

    \item The DeepSeekMath-Base 7B model, which is trained on this dataset, exhibits performance on par with Minerva 540B \citep{minerva}. This outcome suggests that factors other than model size, particularly data quality, significantly influence mathematical reasoning proficiency; a comparatively compact model with superior training data can deliver highly competitive results.

    \item We report insights from a series of math-centric training experiments. Pre-training models on code before engaging in mathematical pre-training measurably enhances their capability to tackle mathematical problems, both in standard reasoning scenarios and when employing external tools. This partially addresses the debated issue of whether pre-training on code augments reasoning skills, especially for mathematics—our results indicate it does.

    \item In contrast to prevailing practices, incorporating arXiv paper data—frequently utilized in many mathematics-focused models—fails to yield observable gains on the full suite of mathematical benchmarks examined in this study.
\end{itemize}

\textbf{Exploration and Analysis of Reinforcement Learning}

\begin{itemize}[topsep=0pt]
    \item We present Group Relative Policy Optimization (GRPO), a resource-efficient and high-performing reinforcement learning method. Unlike traditional approaches such as Proximal Policy Optimization (PPO), GRPO dispenses with the need for a critic model by estimating the baseline directly from group scores. This modification leads to a substantial decrease in training resource requirements.
    \item Through our experiments, we show that GRPO markedly improves the capabilities of our instruction-tuned system, DeepSeekMath-Instruct, using solely instruction-tuning datasets. Additionally, during reinforcement learning, we note notable gains in generalization to out-of-domain tasks.
    \item We propose a consolidated framework that encompasses methods such as RFT, DPO, PPO, and GRPO. Our work also includes extensive empirical studies—such as comparisons between online and offline training, supervision by outcome versus by process, and single-turn against iterative reinforcement learning—to thoroughly examine the core aspects of this unified framework.
    \item Guided by this consolidated framework, we analyze the contributing factors to reinforcement learning's efficacy and outline several prospective avenues to further enhance the reinforcement learning capabilities of large language models.
\end{itemize}

\subsection{Summary of Evaluations and Metrics}

\begin{itemize}[topsep=0pt]
    \item \textbf{English and Chinese Mathematical Reasoning}: We rigorously evaluate our models using diverse English and Chinese benchmark datasets, which span mathematical questions from elementary to advanced undergraduate levels. The English test sets include GSM8K \citep{gsm8k}, MATH \citep{MATH}, SAT \citep{llemma}, OCW Courses \citep{minerva}, and MMLU-STEM \citep{mmlu}, while the Chinese ones feature MGSM-zh \citep{mgsm}, CMATH \citep{wei2023cmath}, Gaokao-MathCloze \citep{agieval}, and Gaokao-MathQA \citep{agieval}. Our evaluations consider both the models’ abilities to generate complete, tool-free textual solutions and their effectiveness in problem-solving via Python code.
    
    For English-language tasks, DeepSeekMath-Base performs on par with the proprietary Minerva 540B \citep{minerva}, and it outperforms every open-source base model available—including Mistral 7B \citep{mistral} and Llemma-34B \citep{llemma}—across both math pre-trained and non-pre-trained models, frequently with considerable advantages. Importantly, DeepSeekMath-Base demonstrates outstanding results on Chinese evaluations as well. This can be attributed to our deviation from earlier strategies \citep{minerva,llemma}, in which pre-training was restricted to English mathematics data; instead, we incorporate high-quality multilingual datasets. Following further instruction tuning and reinforcement learning, DeepSeekMath-Instruct and DeepSeekMath-RL achieve strong outcomes, surpassing the 50\% accuracy threshold for the competitive MATH benchmark—a first for any open-source model.

\item \textbf{Formal Mathematics}: DeepSeekMath-Base is assessed on the informal-to-formal theorem proving benchmark introduced by \citep{dsp_proof}, utilizing miniF2F \citep{minif2f} and selecting Isabelle \citep{isabelle} as the proof assistant. In this setting, DeepSeekMath-Base achieves notable few-shot performance in autoformalization.
\item \textbf{Natural Language Understanding, Reasoning, and Code}: To obtain a well-rounded assessment of the models’ abilities in general understanding, reasoning, and coding, we test DeepSeekMath-Base on several benchmarks. These include the Massive Multitask Language Understanding (MMLU) benchmark \citep{mmlu} covering 57 multiple-choice problems from diverse fields, BIG-Bench Hard (BBH) \citep{bbh} with 23 tasks requiring predominantly multi-step reasoning, as well as HumanEval \citep{codex} and MBPP \citep{mbpp}—standard benchmarks for evaluating code generation models. Results indicate that pre-training on mathematical data enhances both language understanding and reasoning capabilities.

\section{Math Pre-Training}

\subsection{Data Collection and Decontamination} 

Here, we detail the construction of the DeepSeekMath Corpus from Common Crawl data. Figure \ref{fig:data_collect} illustrates our iterative methodology for curating a large-scale mathematical text corpus from Common Crawl, beginning with a carefully selected seed dataset of math-focused texts. This iterative framework is adaptable to domains such as programming as well.

As a starting point, we employ OpenWebMath \citep{openwebmath}, a dataset comprising curated, high-quality mathematical texts from the web, as the initial seed corpus. We leverage this seed to train a fastText classifier \citep{joulin2016fasttext} designed to identify and retrieve additional web pages with similar mathematical content. For training, we randomly sample 500,000 items from the seed as positive instances and select 500,000 web pages from Common Crawl to serve as negatives. Utilizing the fastText open-source implementation\footnote{\url{https://fasttext.cc}}, we set the embedding vector size to 256, use a learning rate of 0.1, allow word n-grams up to length 3, require at least 3 occurrences per word, and conduct 3 training epochs. To reduce dataset redundancy, we apply both URL-level and near-duplicate filtering to Common Crawl, producing a deduplicated subset of 40B HTML pages. Next, the trained fastText model retrieves web pages with high mathematical relevance from this deduplicated corpus. To further enhance quality, retrieved web pages are ranked by the fastText model’s output scores, and only those with the highest scores are retained. The quantity of preserved data is evaluated via pre-training runs at multiple token sizes (40B, 80B, 120B, and 160B tokens), with the first iteration opting for the top 40B tokens.

Following the initial round of data gathering, a significant number of mathematical web pages are still missing from the collected dataset. This is primarily due to the fastText model being trained on a set of positive instances that does not capture a broad enough spectrum of mathematical content. To address this limitation, we expand the seed corpus by sourcing further mathematical websites, enhancing the training data for the fastText model. Our method involves segmenting the entire Common Crawl dataset into non-overlapping domains, each defined as a collection of web pages with the same base URL. For every domain, we determine the proportion of pages captured during the first data collection round. Domains for which more than 10\% of the pages were collected are tagged as math-related (for example, \url{mathoverflow.net}). Next, we conduct manual annotation of URLs pointing to mathematical resources within these designated domains (such as \url{mathoverflow.net/questions}). Any web pages corresponding to these URLs that were missed previously are incorporated into the seed corpus. This process allows us to augment the positive example set, yielding a refined fastText model that can identify a larger volume of mathematical content in subsequent rounds. Across four rounds of data collection, we amass a total of 35.5M mathematical web pages, comprising 120B tokens. By the fourth round, nearly 98\% of the gathered content overlaps with the third round, leading us to conclude the data collection process.

To mitigate the risk of benchmark data leakage, we adopt the approach outlined in \cite{deepseek-coder} to exclude web pages that contain problems or solutions from established English mathematical benchmarks like GSM8K \citep{gsm8k} and MATH \citep{MATH}, as well as Chinese benchmarks such as CMATH \citep{wei2023cmath} and AGIEval \citep{agieval}. Our filtration protocol removes any text block with a 10-gram exact match to any sub-string within these evaluation benchmarks. For benchmark materials under 10 grams but at least 3 grams in length, we use exact matching to identify and eliminate any overlapping pages from our training dataset.

\subsection{Quality Assessment of the DeepSeekMath~Corpus}

We conduct pre-training studies to compare the DeepSeekMath~Corpus against several recently introduced math-focused corpora:

\begin{itemize}[topsep=0pt]
    \item \textbf{MathPile} \citep{mathpile}: a diverse collection containing 8.9B tokens from sources including textbooks, Wikipedia, ProofWiki, CommonCrawl, StackExchange, and arXiv, with more than 85\% of the material originating from arXiv;
    \item \textbf{OpenWebMath} \citep{openwebmath}: a subset of CommonCrawl curated for mathematical content, with a total size of 13.6B tokens;
    \item \textbf{Proof-Pile-2} \citep{llemma}: composed of OpenWebMath, AlgebraicStack (which supplies 10.3B tokens of mathematical code), and arXiv (contributing 28.0B tokens). For experiments on Proof-Pile-2, we adhere to the arXiv:Web:Code ratio of 2:4:1 as described in \cite{llemma}.
\end{itemize}

\subsubsection{Training Setting}
\label{sec:quality-policy}
We perform specialized math training on a pre-existing large language model with 1.3B parameters, utilizing the architecture implemented in DeepSeek LLMs \citep{deepseek-llm}—referred to as DeepSeek-LLM 1.3B. Separate training runs are conducted for each mathematical dataset, each consuming 150B tokens. All training procedures utilize the HAI-LLM \citep{haillm} lightweight and efficient framework. In accordance with the training strategy outlined for DeepSeek LLMs, we employ the AdamW optimizer \citep{adamW} configured with $\beta_1=0.9$, $\beta_2=0.95$, and $\mathrm{weight\_decay}=0.1$, coupled with a multi-phase learning rate scheduler. The learning rate peaks after 2,000 warmup steps, declines to 31.6\% of its maximum by the 80\% completion mark of training, and then drops further to 10.0\% by 90\%. The upper bound for the learning rate is set at $5.3 \times 10^{-4}$, the batch size is maintained at 4M tokens, and the context window is fixed at 4K tokens.

\subsubsection{Evaluation Results}

\textbf{The DeepSeekMath~Corpus stands out for its exceptional quality, extensive multilingual coverage, and unmatched scale among mathematical datasets.}

\begin{itemize}[topsep=0pt]
    \item \textbf{High-quality}:
    We assess the ability of models fine-tuned on different corpora by measuring their performance on eight mathematical benchmarks using few-shot chain-of-thought prompting \cite{cot}. Results presented in Table \ref{tab:corpora_comparison} demonstrate that models trained with the DeepSeekMath~Corpus consistently achieve superior scores. Additionally, Figure \ref{fig:corpora_comparisons} indicates that DeepSeekMath-trained models surpass those trained with Proof-Pile-2 after processing 50B tokens (equivalent to a single epoch for Proof-Pile-2), suggesting that the DeepSeekMath~Corpus offers a higher average data quality.
    \item \textbf{Multilingual}:
    Data within the DeepSeekMath~Corpus spans a wide array of languages, with English and Chinese being the most prevalent. As reflected in Table \ref{tab:corpora_comparison}, models trained on DeepSeekMath~Corpus demonstrate marked improvements in mathematical reasoning across both English and Chinese. By contrast, most existing mathematical corpora are predominantly English-focused, contributing minimal, if any, gains for Chinese-language mathematical reasoning tasks, and may even adversely affect such performance.
    \item \textbf{Large-scale}:
    The DeepSeekMath~Corpus is considerably larger than all previously available mathematical datasets. Figure \ref{fig:corpora_comparisons} illustrates that when DeepSeek-LLM 1.3B is trained with DeepSeekMath~Corpus, the model exhibits a significantly faster and more sustained performance trajectory. Meanwhile, comparative baselines are limited in size, having been cycled through multiple epochs during training, which causes their associated model performance to quickly plateau.
\end{itemize}

\subsection{Training and Evaluating DeepSeekMath-Base 7B}

This section details the development of DeepSeekMath-Base 7B, a foundational model designed to exhibit advanced mathematical reasoning capabilities. The initial model weights are sourced from DeepSeek-Coder-Base-v1.5 7B \citep{deepseek-coder}, followed by training on a corpus totaling 500B tokens. The training data is apportioned as follows: 56\% is drawn from the DeepSeekMath Corpus, 4\% originates from AlgebraicStack, 10\% comprises arXiv content, 20\% consists of Github source code, and the final 10\% includes natural language text sampled from Common Crawl in both English and Chinese. Our approach adheres primarily to the training methodology outlined in Section \ref{sec:quality-policy}, with the exception that the learning rate is capped at 4.2e-4, and each batch includes 10M tokens.

To thoroughly evaluate the mathematical proficiency of DeepSeekMath-Base 7B, we examine its ability to independently generate complete mathematical solutions, utilize external tools to solve problems, and perform formalized theorem proving. In addition to its mathematical strength, we present broader evaluations, highlighting the model's competencies in natural language understanding, deductive reasoning, and coding.

\paragraph{Step-by-Step Reasoning in Mathematical Problem Solving}

DeepSeekMath-Base is assessed on its stepwise reasoning and problem-solving abilities through few-shot chain-of-thought prompting \citep{cot} on a suite of eight benchmarks in both English and Chinese. These evaluations incorporate a range of quantitative reasoning tasks—such as those from GSM8K \citep{gsm8k}, MATH \citep{MATH}, and CMATH \citep{wei2023cmath}—as well as multiple-choice examinations including MMLU-STEM \citep{mmlu} and Gaokao-MathQA \citep{agieval}. The selected benchmarks collectively span a comprehensive range of mathematical domains from primary through collegiate levels.

As detailed in Table \ref{tab:base_math_cot}, DeepSeekMath-Base 7B achieves the highest scores across all eight evaluated benchmarks among publicly available base models, such as the extensively used Mistral 7B \citep{mistral} and the newly introduced Llemma 34B \citep{llemma}, both trained with mathematical corpora including Proof-Pile-2 \citep{llemma}. In particular, on the advanced MATH dataset, DeepSeekMath-Base not only exceeds other open-source base models by more than 10\% in absolute terms, but also surpasses Minerva 540B \citep{minerva}, a proprietary model that is 77 times larger and incorporates additional training on mathematical materials using PaLM \citep{palm}.

\paragraph{Mathematical Problem Solving with Tool Use} To assess the models’ capacity for program-assisted mathematical reasoning, we apply few-shot program-of-thought prompting \citep{pot,pal} to GSM8K and MATH. For each problem, the models are instructed to compose a Python script—potentially employing libraries such as \textit{math} and \textit{sympy}—to carry out complex calculations. The output of this code serves as the predicted solution. As indicated by the results in Table \ref{tab:base_math_pot_proof}, DeepSeekMath-Base 7B achieves superior results compared to the previous leading model, Llemma 34B.

\paragraph{Formal Mathematics}

Automated formal proof generation is increasingly regarded as crucial for improving both the correctness and efficiency of mathematical proofs. We assess DeepSeekMath-Base 7B on the informal-to-formal proof translation task presented in \citep{dsp_proof}, which involves generating a formal proof using an informal claim, its formal translation, and an informal proof sketch. Evaluations take place on miniF2F \citep{minif2f}, which benchmarks formal reasoning at Olympiad difficulty, with each instance requiring a formal proof in Isabelle provided via few-shot prompting. In line with \cite{dsp_proof}, the approach involves model generation of proof outlines and subsequent completion of these sketches using the automated theorem prover Sledgehammer \citep{sledgehammer}. As displayed in Table \ref{tab:base_math_pot_proof}, DeepSeekMath-Base 7B achieves robust performance in the task of autoformalizing mathematical proofs.

\paragraph{Natural Language Understanding, Reasoning, and Code}

Model capabilities for natural language comprehension are evaluated on MMLU \citep{mmlu}, for reasoning on BBH \citep{bbh}, and for coding abilities using HumanEval \citep{codex} and MBPP \citep{mbpp}. Results in Table \ref{tab:understanding_reasoning_code} demonstrate that DeepSeekMath-Base 7B achieves considerable performance gains on MMLU and BBH compared to the prior DeepSeek-Coder-Base-v1.5 \citep{deepseek-coder}, signifying that mathematics-focused training enhances general language and reasoning skills. Furthermore, by integrating code tokens into continued training, DeepSeekMath-Base 7B sustains the code generation performance of DeepSeek-Coder-Base-v1.5 on the programming benchmarks. Overall, DeepSeekMath-Base 7B consistently exceeds the performance of the general-purpose Mistral 7B \citep{mistral} on all three benchmarks pertaining to reasoning and code.

\section{Supervised Fine-Tuning}

\subsection{SFT Data Curation}

We assembled a comprehensive instruction-tuning dataset for mathematics that encompasses both English and Chinese problems. The dataset spans a range of mathematical domains and features tasks of varying complexity; each problem is coupled with a solution presented in chain-of-thought (CoT) \citep{cot}, program-of-thought (PoT) \citep{pot,pal}, or a tool-integrated reasoning format \citep{tora}. Altogether, our training set consists of 776K examples.

\begin{itemize}[topsep=0pt]
    \item \textbf{English mathematical datasets}: Our English-language set features GSM8K and MATH, which we enrich by annotating their problems with tool-integrated solutions. We further supplement this set using selected items from MathInstruct \citep{MathInstruct} and the training partition of Lila-OOD \citep{lila}, in which tasks are solved with either CoT or PoT methodologies. The resulting collection spans fields such as algebra, probability, number theory, calculus, and geometry.
    \item \textbf{Chinese mathematical datasets}: Our Chinese dataset draws on K-12 mathematics problems covering 76 distinct subtopics—including, for example, linear equations—with answers provided in both CoT and tool-integrated reasoning formats.
\end{itemize}

\subsection{Training and Evaluating DeepSeekMath-Instruct 7B}

We present DeepSeekMath-Instruct 7B, a model that receives instruction tuning in mathematics atop the DeepSeekMath-Base architecture. During training, we concatenate randomly selected samples until reaching the 4K token context length constraint. The model is trained for 500 steps, employing a batch size of 256 and maintaining a fixed learning rate of 5e-5.

For evaluation, we assess the mathematical capabilities of the models under two regimes: with and without tool integration. Testing is performed on four quantitative reasoning benchmarks encompassing both English and Chinese. Comparisons are drawn against state-of-the-art contemporaries:

\begin{itemize}[topsep=0pt]
    \item \textbf{Closed-source models} include: (1) various members of the GPT family, notably GPT-4 \citep{gpt4} and GPT-4 Code Interpreter~\footnote{\url{https://openai.com/blog/chatgpt-plugins\#\#code-interpreter}}, (2) Gemini Ultra and Pro \citep{gemini}, (3) Inflection-2 \citep{inflection-2}, (4) Grok-1~\footnote{\url{https://x.ai/model-card}}, and additional recent entrants from Chinese organizations such as (5) Baichuan-3~\footnote{\url{https://www.baichuan-ai.com}}, (6) the latest GLM-4~\footnote{\url{https://open.bigmodel.cn/dev/api\#glm-4}}

    from the GLM family \citep{glm}.
\end{itemize}

These models are designed for a wide range of applications and have generally been subjected to various alignment processes. 
\item \textbf{Open-source models} encompass: general-purpose systems such as (1) DeepSeek-LLM-Chat 67B \citep{deepseek-llm}, (2) Qwen 72B \citep{qwen}, (3) SeaLLM-v2 7B \citep{seallm}, and (4) ChatGLM3 6B \citep{chatglm3}; as well as those models that have been specifically tailored for enhanced mathematical abilities, including (5) InternLM2-Math 20B~\footnote{\url{https://github.com/InternLM/InternLM-Math}}, an extension of InternLM2 further trained on mathematical data and instruction-tuned; (6) Math-Shepherd-Mistral 7B, which applies PPO \citep{schulman2017proximal} to the Mistral 7B \citep{mistral} model using a process-supervised reward framework; (7) the WizardMath family \citep{wizardmath}, which advances math reasoning in Mistral 7B and Llama-2 70B \citep{llama2} through the use of evolve-instruct—a variant of instruction tuning with AI-generated instructions—and PPO on datasets predominantly composed of GSM8K and MATH tasks; (8) MetaMath 70B \citep{metamath}, achieved by fine-tuning Llama-2 70B on an augmented GSM8K and MATH dataset; (9) ToRA 34B \cite{tora}, in which CodeLlama 34B is further trained to execute mathematical reasoning augmented with external tools; and (10) MAmmoTH 70B \citep{MathInstruct}, where Llama-2 70B receives instruction tuning on MathInstruct.

Table \ref{tab:sft_rl_math} illustrates that, in conditions where external tools are not permitted, DeepSeekMath-Instruct 7B achieves robust stepwise mathematical reasoning. Particularly, this model achieves the highest score on the MATH benchmark among open-source solutions and even exceeds the majority of proprietary models, such as Inflection-2 and Gemini Pro, by a margin of at least 9 percentage points, even when compared to models that are significantly larger (e.g., Qwen 72B) or those specifically improved by reinforcement learning on mathematical data (e.g., WizardMath-v1.1 7B). DeepSeekMath-Instruct competes closely with proprietary Chinese models such as GLM-4 and Baichuan-3 on the MATH benchmark, although it remains outperformed by GPT-4 and Gemini Ultra.

In a different evaluation context that allows for both natural language reasoning and the utilization of programmatic tools to solve problems, DeepSeekMath-Instruct 7B attains an accuracy approaching 60\% on the MATH dataset, outperforming all open-source alternatives currently available. On other evaluated benchmarks, our model achieves results on par with DeepSeek-LLM-Chat 67B, the previous open-source leader, despite being only one-tenth its size.

\section{Reinforcement Learning}

\subsection{Group Relative Policy Optimization} Reinforcement learning (RL) has been shown to effectively boost the mathematical reasoning performance of LLMs subsequent to the Supervised Fine-Tuning (SFT) phase \citep{wang2023math,wizardmath}. Here, we present our novel RL method, Group Relative Policy Optimization (GRPO), designed for improved efficiency and effectiveness.

\subsubsection{From PPO to GRPO} Proximal Policy Optimization (PPO) \citep{schulman2017proximal} stands out as an actor-critic reinforcement learning method that is prevalently used in RL fine-tuning for LLMs \citep{ouyang2022training}. Its optimization process targets maximizing the following surrogate function:

\begin{equation}
\footnotesize
    \mathcal{J}_{PPO}(\theta) = \mathbb{E}{[q \sim P(Q), o \sim \pi_{\theta_{old}}(O|q)]} \frac{1}{|o|} \sum_{t=1}^{|o|} \min \left[ \frac{\pi_\theta(o_{t} | q, o_{<t})}{\pi_{\theta_{old}}(o_{t} | q, o_{<t})} A_{t}, \text{clip} \left( \frac{\pi_\theta(o_{t} | q, o_{<t})}{\pi_{\theta_{old}}(o_{t} | q,

Here, $\pi_{\theta}$ and $\pi_{\theta_{old}}$ denote the present and previous policy models, respectively. The variables $q$ and $o$ represent questions and outputs, sampled from both the question dataset and from $\pi_{\theta_{old}}$. The hyper-parameter $\varepsilon$ governs the extent of clipping in PPO to enhance training stability. The term $A_t$ stands for the advantage estimate, derived using Generalized Advantage Estimation (GAE) \citep{gae}, which relies on reward values $\{r_{\ge t}\}$ and a learned value function $V_{\psi}$. As such, PPO requires concurrent training of a value function alongside the policy model. To prevent excessive optimization against the reward model, it is common practice to incorporate a per-token KL divergence penalty between the learned policy and a fixed reference model at each token in the reward function \citep{ouyang2022training}, as follows:

\begin{equation}
    r_{t} = r_\varphi(q, o_{\le t}) - \beta \log\frac{\pi_{\theta}(o_{t}|q, o_{<t})}{\pi_{ref}(o_{t}|q, o_{<t})},
\label{eq:PPO-reward}
\end{equation}

where $r_\varphi$ is the reward function, $\pi_{ref}$ refers to the reference policy (often the original SFT model), and $\beta$ is the KL regularization coefficient.

Because the PPO framework typically employs a value network as large as the main policy network, the resulting memory and computation costs are significant. In addition, PPO leverages the value function as a variance-reducing baseline for the advantage computation. However, in large language models, the reward model frequently assigns a nonzero score only to the last token, posing challenges for learning a precise per-token value estimate. To overcome these issues, we introduce Group Relative Policy Optimization (GRPO) as illustrated in Figure \ref{fig:grpo}. GRPO eliminates the explicit training of an auxiliary value model by instead utilizing the mean reward from a set of multiple outputs sampled for the same question as a groupwise baseline. In detail, for each question $q$, GRPO samples a collection of outputs $\{o_1, o_2, \cdots, o_G\}$ from $\pi_{\theta_{old}}$, and then updates the policy by maximizing:

\begin{equation}
\footnotesize
\begin{split}
    \mathcal{J}_{GRPO}(\theta) &= \mathbb{E}{[q \sim P(Q), \{o_i\}_{i=1}^G \sim \pi_{\theta_{old}}(O|q)]}  \\
    & \frac{1}{G}\sum_{i=1}^G\frac{1}{|o_i|} \sum_{t=1}^{|o_i|} \left\{ \min \left[ \frac{\pi_\theta(o_{i,t} | q, o_{i,<t})}{\pi_{\theta_{old}}(o_{i,t} | q, o_{i,<t})} \hat{A}_{i,t}, \text{clip} \left( \frac{\pi_\theta(o_{i,t} | q, o_{i,<t})}{\pi_{\theta_{old}}(o_{i,t} | q, o_{i,<t})}, 1 - \varepsilon, 1 + \varepsilon \right)  \hat{A}_{i,t} \right] - \beta \mathbb{D}_{KL}\left[\pi_{\theta} || \pi_{ref}\right]\right\} ,
\end{split}
\label{eq:GRPO-obj}
\end{equation}

where both $\varepsilon$ and $\beta$ function as hyper-parameters, and $\hat{A}_{i,t}$ denotes the advantage based exclusively on within-group relative rewards, whose computation is explained in detail in subsequent subsections. The GRPO strategy for constructing advantage estimates aligns with the comparative design of reward models, which are typically trained using paired output comparisons for the same input. Notably, rather than incorporating a KL penalty inside the reward, GRPO introduces the KL divergence between the policy and reference policy directly into the loss objective, thereby streamlining the evaluation of $\hat{A}_{i,t}$. Distinct from the KL term employed in (\ref{eq:PPO-reward}), the KL divergence in this case is computed using the following unbiased estimator \citep{kl_approx}:

\begin{equation}
\small
    \mathbb{D}_{KL}\left[\pi_{\theta} || \pi_{

\begin{algorithm}[t]
  \small
  \caption{Iterative Group Relative Policy Optimization}
  \textbf{Input} initial policy model $\pi_{\theta_{\text{init}}}$; reward models $r_\varphi$; task prompts $\mathcal{D}$; 
  hyperparameters $\varepsilon$, $\beta$, $\mu$
  \begin{algorithmic}[1]
    \State Set policy model $\pi_\theta$ to $\pi_{\theta_{\text{init}}}$
    \For{iteration = 1, \dots, I}
       \State Assign reference model $\pi_{ref}$ to current policy $\pi_{\theta}$
      \For{step = 1, \dots, M}
      \State Select a mini-batch $\mathcal{D}_b$ from $\mathcal{D}$
      \State Update the previous policy model $\pi_{\theta_{old}}$ with current parameters $\pi_{\theta}$ 
      \State For each question $q \in \mathcal{D}_b$, generate $G$ outputs $\{o_i\}_{i=1}^G$ using $\pi_{\theta_{old}} (\cdot \mid q)$ 
      \State Evaluate each output $o_i$ with $r_{\varphi}$, resulting in rewards $\{r_i\}_{i=1}^{G}$
      \State For each $t$-th token of $o_i$, compute $\hat{A}_{i,t}$ using group-relative advantage estimation
      \For{GRPO iteration = 1, \dots, $\mu$}
        \State Adjust policy model $\pi_{\theta}$ by optimizing the GRPO loss function (Equation \ref{eq:GC-GRPO})
      \EndFor
    \EndFor 
    \State Continuously update $r_\varphi$ with replay-based training
    \EndFor 
  \end{algorithmic}
  \textbf{Output} $\pi_\theta$
  \label{alg:iter-grpo}
\end{algorithm}

\subsubsection{Outcome Supervision RL with GRPO} More formally, for each question $q$, $G$ output samples $\{o_1, o_2, \cdots, o_G\}$ are drawn from the prior policy $\pi_{\theta_{old}}$. These outputs are evaluated by the reward model to obtain a set of corresponding reward values $\mathbf{r}=\{r_1, r_2, \cdots, r_G\}$. The rewards are then standardized by subtracting their mean and dividing by their standard deviation across the group. Under outcome supervision, this standardized reward is provided once at the terminal token of each $o_i$, and all token-level advantages $\hat{A}_{i, t}$ throughout $o_i$ are assigned this value, specifically, $\hat{A}_{i, t} = \widetilde{r}_i = \frac{r_i- {\rm mean}(\mathbf{r})}{{\rm std}(\mathbf{r})}$. Policy optimization then seeks to maximize the objective given in equation (\ref{eq:GRPO-obj}).

\subsubsection{Process Supervision RL with GRPO} As outcome supervision delivers feedback only at completion, it may be inadequate for effectively steering policies on complex mathematical problems. Building on \cite{wang2023math}, we additionally investigate process supervision, where feedback is provided after each intermediate reasoning step. Explicitly, for a prompt $q$ and $G$ sampled completions $\{o_1, o_2, \cdots, o_G\}$, a process-oriented reward model assigns scores to each reasoning step, resulting in reward sets $\mathbf{R} = \{ \{r_1^{index(1)},\cdots,r_1^{index(K_1)}\}, \cdots,  \{r_G^{index(1)},\cdots,r_G^{index(K_G)}\} \}$, where $index(j)$ indicates the position of the $j$-th step’s final token, and $K_i$ represents the step count in output $i$. These stepwise rewards are also standardized by the mean and standard deviation, yielding $\widetilde{r}_i^{index(j)} = \frac{r_i^{index(j)} - {\rm mean(\mathbf{R})}}{{\rm std(\mathbf{R})}}$. 
Process supervision then defines each token’s advantage as the sum of normalized future step rewards: $\hat{A}_{i, t} = \sum_{index(j) \ge t} \widetilde{r}_i^{index(j)}$. The resulting policy is improved by maximizing the target from equation (\ref{eq:

\subsubsection{Iterative RL with GRPO} During reinforcement learning, the reward model used at earlier stages may become inadequate for effectively guiding the policy model as training advances. To address this, we investigate an iterative RL strategy utilizing GRPO. Algorithm \ref{alg:iter-grpo} details this process: new reward model training datasets are periodically created from samples generated by the latest policy model. The reward model is updated via a replay buffer that includes 10\% of past data. Subsequently, the reference model is synchronized with the current policy model, after which the policy model is updated further using feedback from the refreshed reward model.

\subsection{Training and Evaluating DeepSeekMath-RL}

Reinforcement learning is applied on top of DeepSeekMath-Instruct 7B. The RL training corpus is drawn from chain-of-thought-style problems associated with GSM8K and MATH, extracted from the SFT dataset and totaling approximately 144K items. Other SFT-derived questions are excluded in order to specifically evaluate the effect of RL on benchmarks that are underrepresented during RL training. Construction of the reward model training data mirrors the approach in \citep{wang2023math}. The initial reward model is trained on DeepSeekMath-Base 7B with a learning rate of 2e-5. For GRPO, the policy model adopts a learning rate of 1e-6, and the KL penalty coefficient is set at 0.04. We generate $64$ sampled responses for each question, with a maximum sequence length of 1024 tokens and a batch size of 1024 for training. Policy updates occur just once following each exploration round. For evaluation, DeepSeekMath-RL 7B is benchmarked following the same setup as DeepSeekMath-Instruct 7B, where GSM8K and MATH tasks involving chain-of-thought are categorized as in-domain, while the remaining benchmarks are considered out-of-domain.

Table \ref{tab:sft_rl_math} presents a comparative analysis of open- and closed-source models utilizing both chain-of-thought and tool-integrated reasoning across English and Chinese benchmarks. Our observations are as follows: 1) With chain-of-thought prompting, DeepSeekMath-RL 7B achieves accuracy rates of 88.2\% on GSM8K and 51.7\% on MATH. These results exceed the performance of all open-source models within the 7B to 70B range and outperform most closed-source counterparts. 2) Importantly, the DeepSeekMath-RL 7B model’s instruction tuning phase is limited solely to chain-of-thought-style data from GSM8K and MATH, initializing from DeepSeekMath-Instruct 7B. Despite the restricted diversity of its training corpus, DeepSeekMath-RL 7B surpasses DeepSeekMath-Instruct 7B on all metrics, thus highlighting the notable impact of reinforcement learning.

\section{Discussion}
This section conveys insights gained from our investigations into pre-training and reinforcement learning.

\subsection{Lessons Learnt in Pre-Training}

We begin by outlining lessons from our pre-training process. Unless mentioned otherwise, training configurations follow those specified in Section \ref{sec:quality-policy}. Additionally, references to the DeepSeekMath Corpus in this section pertain to the 89B-token dataset assembled during the second round of data gathering.

\subsubsection{Code Training Benefits Mathematical Reasoning}

There exists a widely held but not conclusively validated supposition that training with code enhances reasoning capabilities. Here, we provide empirical evidence within the mathematical context: incorporating code training enhances a model’s mathematical reasoning, with or without auxiliary tools.

To evaluate the influence of code-focused training on mathematical reasoning abilities, we tested two approaches: a two-stage training pipeline and a single-stage training baseline.

\textbf{Two-Stage Training}

\begin{itemize}[topsep=0pt]
    \item \textbf{Code Training for 400B Tokens $\rightarrow$ Math Training for 150B Tokens}: DeepSeek-LLM 1.3B is first exposed to 400B code tokens and then continues training on an additional 150B math tokens.
    \item \textbf{General Training for 400B Tokens $\rightarrow$ Math Training for 150B Tokens}: As a baseline comparison, we substitute the code tokens in the initial phase with general tokens (sourced from an extensive general-purpose corpus curated by DeepSeek-AI) to evaluate the differential impact of code versus general data on mathematical reasoning enhancement.
\end{itemize}

\textbf{One-Stage Training}

\begin{itemize}[topsep=0pt]
    \item \textbf{Math Training for 150B Tokens}: DeepSeek-LLM 1.3B undergoes training exclusively with 150B math tokens.
    \item \textbf{Training on a mixture of 400B Code Tokens and 150B Math Tokens}: Since sequential math training post-code training can diminish coding capabilities, we also analyze whether co-training on a combined dataset (400B code tokens + 150B math tokens) in a single phase can bolster mathematical reasoning and counteract catastrophic forgetting.
\end{itemize}

\paragraph{Results}
Performance across these experimental regimes is presented in Table \ref{tab:code-math} and Table \ref{tab:code-math-general-eval}.

Incorporating code training demonstrably enhances program-based mathematical reasoning under both the two-stage and one-stage paradigms. According to Table \ref{tab:code-math}, code pretraining in a two-phase setup already markedly increases proficiency on tasks such as GSM8K and MATH when solved via Python, with the subsequent math-focused training providing additional gains. Notably, in the one-stage framework, jointly training on code and math data addresses the catastrophic forgetting observed in sequential training and further augments both coding aptitude (Table \ref{tab:code-math-general-eval}) and program-supported mathematical reasoning (Table \ref{tab:code-math}).

Code-focused pretraining also advances mathematical reasoning in settings without external tools. Within the two-phase schema, an initial pass on code data imparts meaningful improvements and amplifies the effect of later math training, resulting in peak performance. However, in the single-stage training where code and math data are blended, a drop in tool-free mathematical reasoning ability is observed. One possible explanation is that the 1.3B parameter scale of DeepSeek-LLM limits its capacity to simultaneously master both code and math information when trained in this joint manner.

\subsubsection{ArXiv Papers Appear Unhelpful for Enhancing Mathematical Reasoning}

While arXiv papers are frequently incorporated into the pre-training datasets for mathematical language models \citep{minerva,gpt-f,llemma,mathpile}, thorough investigations into their contribution to mathematical reasoning remain scarce. Interestingly, our experimental results suggest that arXiv papers do not contribute significantly to advances in mathematical reasoning capability. We conducted assessments across multiple model architectures and sizes, specifically DeepSeek-LLM 1.3B and DeepSeek-Coder-Base-v1.5 7B \citep{deepseek-coder}, drawing on arXiv-based corpora that have undergone different data preparation processes:

\begin{itemize}[topsep=0pt]
    \item \textbf{MathPile} \citep{mathpile}: a dataset comprising 8.9 billion tokens curated using heuristic filtering and cleaning strategies, with scientific arXiv content making up more than 85\%;
    \item \textbf{ArXiv-RedPajama} \citep{redpajama}: a 28.0 billion-token dataset consisting of arXiv LaTeX documents, stripped of elements such as preambles, comments, macros, and references.
\end{itemize}

For our analysis, DeepSeek-LLM 1.3B was individually trained on 150B tokens and DeepSeek-Coder-Base-v1.5 7B on 40B tokens from each arXiv dataset. The findings indicate a lack of measurable benefit from arXiv corpora with respect to mathematical reasoning. When models are trained exclusively with arXiv-derived data, no clear gains—and sometimes regressions—are observed across a range of mathematical benchmarks with varying levels of complexity employed in this work. Evaluation suites encompass quantitative reasoning sets such as GSM8K and MATH (see Table \ref{tab:arxiv-cot}), multiple-choice assessments like MMLU-STEM (see Table \ref{tab:arxiv-cot}), as well as formal mathematics challenges including miniF2F (see Table \ref{tab:arxiv_atp}).

Nonetheless, these findings are subject to several caveats, and definitive judgments should be avoided. Open questions include:

\begin{itemize}[topsep=0pt]
    \item Possible contributions of arXiv-derived data to particular mathematical tasks not addressed herein, for example, transforming formal statements or proofs into informal text;
    \item The interaction effects when arXiv tokens are mixed with data of other types;
    \item Whether scaling models to larger sizes might reveal latent advantages of arXiv papers.
\end{itemize}

Consequently, these issues necessitate deeper investigation, which is reserved for subsequent research efforts.

\subsection{Perspectives on Reinforcement Learning}
\subsubsection{Towards a Unified Framework} This section introduces a generalized framework for comparing various training algorithms, including SFT, RFT, DPO, PPO, GRPO, and proceeds to empirically examine the elements constituting this framework. Typically, for these approaches, the gradient with respect to parameters $\theta$ is formalized as follows:

\begin{equation}
    \nabla_{\theta}\mathcal{J}_{\textcolor{red}{\mathcal{A}}}(\theta) = \mathbb{E}[\underbrace{(q,o) \sim \textcolor{red}{\mathcal{D}}}_{Data \ Source}]\left( \frac{1}{|o|} \sum_{t=1}^{|o|}  \underbrace{GC_{{\mathcal{A}}}(q, o, t, \textcolor{red}{\pi_{{rf}}})}_{Gradient \ Coefficient}  \nabla_{\theta}\log \pi_{\theta}(o_t | q, o_{<t})\right).
\label{eq:objective}
\end{equation}

This framework is characterized by three principal elements: 1) the \textit{Data Source $\mathcal{D}$}, specifying the dataset utilized during training; 2) the \textit{Reward Function $\pi_{{rf}}$}, which generates the reward signals guiding model optimization; and 3) the \textit{Algorithm $\mathcal{A}$}, which translates training data and reward signals into the gradient coefficient $GC$, dictating the extent of model updates via penalty or reinforcement. Several prominent training approaches can be interpreted through this unified lens:

\begin{itemize}
    \item  \textbf{Supervised Fine-tuning (SFT)}: SFT involves adapting a pretrained model with human-curated SFT examples.
    \item \textbf{Rejection Sampling Fine-tuning (RFT)}: RFT continues the fine-tuning process using outputs sampled from the SFT model and filtered according to answer correctness, based on SFT prompts.
    \item \textbf{Direct Preference Optimization (DPO)}: DPO enhances the SFT model by applying fine-tuning with pairwise loss on extra sampled outputs generated by the SFT model.
    \item \textbf{Online Rejection Sampling Fine-tuning (Online RFT)}: Distinct from traditional RFT, Online RFT initializes the policy model from the SFT weights, but performs fine-tuning with augmented outputs generated by the current (online) policy.
    \item \textbf{PPO/GRPO}: These algorithms start with the SFT model for initialization, and further reinforce it using outputs sampled in real-time from the policy as it evolves.
\end{itemize}

The constituent elements of these approaches are outlined in Table \ref{tab:data_policy}. For an expanded derivation, please consult Appendix \ref{app:analysis-rl}.

\paragraph{Considerations on Data Source} We categorize data sources into online and offline sampling. Online sampling refers to datasets collected directly from the real-time exploration of the current policy during training, whereas offline sampling makes use of data derived from the behavior of the initial SFT model. RFT and DPO adhere to the offline paradigm, while Online RFT and GRPO implement the online approach.

Figure \ref{fig:rl-analysis} illustrates that Online RFT exhibits marked improvement over RFT on both evaluation benchmarks. In the early phases of training, Online RFT and RFT display similar performance; however, as training progresses, Online RFT achieves a pronounced lead, reflecting the efficacy of online sampling. This is understandable since, at the outset, the actor and SFT model produce similar outputs, with sampled examples exhibiting minimal variance. As training advances, samples from the evolving actor increasingly diverge from those of the SFT model, amplifying the benefits of real-time sampling.

\paragraph{Considerations on Gradient Coefficient} The transformation of the reward signal into a gradient coefficient is pivotal for model updates. In our analysis, we bifurcate the reward function into 'Rule' and 'Model' categories: 'Rule' involves evaluating answer quality by correctness, whereas 'Model' entails leveraging a trained reward model (itself trained on data labeled by the rule-based judgment) to assess responses. Equations \ref{eq:GC-RFT} and \ref{eq:GC-GRPO} emphasize a principal distinction between GRPO and Online RFT: GRPO dynamically calibrates its gradient coefficient contingent upon the reward model's output, allowing it to provide differentiated reinforcement and penalty based on the magnitude of the rewards. Conversely, Online RFT does not perform such differentiation, applying the same degree of positive reinforcement to all correct answers without penalizing incorrect ones.

As illustrated in Figure \ref{fig:rl-analysis}, GRPO achieves superior results compared to online RFT, underscoring the effectiveness of adjusting the positive and negative gradient coefficients. Moreover, GRPO+PS outperforms GRPO+OS, demonstrating the advantage of employing fine-grained, step-aware coefficients during gradient calculation. Additionally, we investigate iterative reinforcement learning, conducting two cycles of iteration in our experiments. Figure \ref{fig:iter-rl} reveals that iterative RL substantially boosts model performance, with the first iteration contributing the most significant gains.

\subsubsection{Why RL Works?}  
Within this study, we apply reinforcement learning on a selected subset of instruction tuning data, yielding marked improvements over the baseline instruction-tuned model. To elucidate the mechanism underlying reinforcement learning's effectiveness, we assess the Pass@K and Maj@K accuracy metrics for both the Instruct and RL-optimized models on two separate benchmarks. As depicted in Figure \ref{fig:maj-pass}, reinforcement learning yields noticeable gains in Maj@K but does not improve Pass@K. These observations suggest that RL primarily fortifies the model’s output distribution, making it more robust; put differently, \textbf{the observed improvements arise from promoting the selection of the correct response from the TopK outputs, as opposed to augmenting the model's underlying capability.} This phenomenon aligns with findings by \citep{wang2023making}, who identified a \textbf{misalignment problem} in reasoning tasks within SFT models, and demonstrated that a sequence of preference alignment strategies \citep{yuan2023rrhf,song2023preference,wang2023making} can enhance the reasoning ability in such settings.

\subsubsection{How to Achieve More Effective RL?}
\label{sec:effec_rl}
Our results validate that RL is highly effective for mathematical reasoning challenges. We further introduce a unified conceptual framework for interpreting a range of prominent training approaches, treating them uniformly as forms of either explicit or approximated RL techniques. According to Equation \ref{eq:objective}, three principal elements define these methods: Data Source, Algorithm, and Reward Function. We discuss potential research avenues for each of these components.

\paragraph{Data Source} The data source underpins all training methodologies. For RL, this typically means unlabeled prompts, with candidate outputs generated by sampling from the policy model. In our experiments, we restrict ourselves to questions from the instruction tuning set and utilize basic nucleus sampling for output generation. We speculate this limitation partially explains why our RL framework primarily enhances Maj@K. In future work, we aim to extend RL to out-of-distribution prompts and integrate \textbf{advanced sampling (decoding) strategies}, such as those utilizing tree-search methodologies \citep{yao2023tree}. Furthermore, adopting \textbf{efficient inference techniques} \citep{xia-etal-2023-speculative,leviathan2023fast,kwon2023efficient,xia2024unlocking}, which are crucial for maximizing the exploration capabilities of policy models, will also be a significant focus.

\paragraph{Algorithms} Algorithms utilize both data and reward signals to compute the gradient coefficient, which in turn guides updates to the model parameters. According to Equation \ref{eq:objective}, existing methods predominantly \textbf{TRUST} the output of the reward function to modulate the conditional probability of specific tokens, either increasing or decreasing it. Nonetheless, the reliability of the reward signal cannot be consistently guaranteed, especially for highly intricate tasks. As an illustration, the PRM800K datasets \citep{lightman2023let}—despite having been meticulously labeled by experienced annotators—still contain around 20\% erroneous annotations\footnote{\url{https://github.com/openai/prm800k/issues/12\#issuecomment-1728491852}}. Therefore, our investigation will focus on reinforcement learning algorithms capable of maintaining robustness in the face of noisy reward signals. We posit that such \textbf{WEAK-TO-STRONG} \citep{burns2023weak} alignment strategies could revolutionize learning paradigms.

\paragraph{Reward Function} The reward function serves as the primary training signal provider. Within RL, this typically manifests as a neural reward model. We identify three pivotal avenues for advancing reward models: 1) \textbf{Improving the generalization capacity of the reward model.} Reward models should effectively generalize to manage both out-of-distribution queries and sophisticated decoding results; failing to do so may lead reinforcement learning to merely reinforce existing distributions within LLMs without genuine advancement of core capabilities; 2) \textbf{Capturing and expressing the uncertainty in reward models.} Representing uncertainty could function as an integrative link between weak reward functions and weak-to-strong learning frameworks; 3) \textbf{Developing efficient methods to construct high-quality process reward models} capable of furnishing nuanced training signals to facilitate complex reasoning processes \citep{lightman2023let,wang2023math}.

\section{Conclusion, Limitation, and Future Work} This work introduces DeepSeekMath, which surpasses all available open-source models on the challenging MATH benchmark and approaches the results demonstrated by proprietary systems. Built upon DeepSeek-Coder-v1.5 7B, DeepSeekMath~is continually trained on 500B tokens, with a noteworthy fraction comprising 120B mathematical tokens harvested from Common Crawl. Through comprehensive ablation experiments, we demonstrate that web-sourced pages are a rich resource for high-caliber mathematical content, whereas data from arXiv may offer less benefit than initially anticipated. We propose Group Relative Policy Optimization (GRPO), an adaptation of Proximal Policy Optimization (PPO), which significantly enhances mathematical reasoning while being more memory-efficient. Empirical findings reveal that GRPO remains beneficial even as DeepSeekMath-Instruct 7B achieves strong results across benchmarks. Additionally, we offer a unifying framework for understanding various reinforcement learning strategies and highlight several promising avenues for the development of more effective reinforcement learning approaches.

Despite its substantial success on quantitative reasoning evaluations, DeepSeekMath~shows comparatively weaker performance in geometry and theorem-proving tasks, relative to proprietary models. For example, in preliminary trials, the model struggled with problems involving geometric figures such as triangles and ellipses, suggesting a possible pre-training and fine-tuning dataset bias. Furthermore, limited by its model size, DeepSeekMath~falls short of GPT-4 in few-shot settings: while GPT-4's performance improves with additional examples, DeepSeekMath~exhibits similar capabilities regardless of zero-shot or few-shot context. Moving forward, we plan to further enhance our curated data selection process to assemble an even higher quality pre-training corpus. Moreover, we will investigate the future directions outlined in Section \ref{sec:effec_rl} for advancing reinforcement learning techniques applicable to LLMs.

\end{document}